% SPDX-License-Identifier: GPL-3.0-or-later
% Copyright (C) 2026 Alessandro Greco
%% EPUB 无障碍替代文本。
\altfig{本书 Chutney 测试网络的所有进程均绑定在 127.0.0.1：三台兼任中继的目录权威节点 000a 至 002a，使用 ORPort 5100–5102 和 DirPort 7100–7102；三台中间中继 003r 至 005r、两台出口中继 006e 和 007e、一台使用 SOCKS 9008 的客户端 008c，以及洋葱服务 009h，其虚拟端口 5858 映射到 127.0.0.1:4747。实验中权威节点每 20 秒投票一次。HTTP 示例电路从客户端依次经过 003r、005r 和出口 006e，抵达本地测试服务器 127.0.0.1:8080；实验出口策略专门允许访问本地地址。洋葱服务是独立的，不经过出口中继。图中仅列部分端口，完整端口表见章节。绑定环回地址不等同于与宿主机网络隔离。}
% 实验室私有网络：简化拓扑
\begin{tikzpicture}[font=\sffamily\scriptsize,
  aut/.style={relay,fill=torviola!15,draw=torviola},
  ex/.style={relay,fill=torverde!30}]
  \foreach \i/\x in {0/0,1/2,2/4} \node[aut] (a\i) at (\x,3.2) {00\i a};
  \node[above=1pt of a1,font=\sffamily\scriptsize,align=center] {目录权威节点\\ORPort 5100--5102, DirPort 7100--7102};
  \node[font=\sffamily\scriptsize,text=torgrigio,align=center] at (6.6,3.2) {每 20 秒投票\\签名共识};
  \foreach \i/\x in {3/0,4/2,5/4} \node[relay] (r\i) at (\x,1.2) {00\i r};
  \node[ex] (e6) at (6.6,1.2) {006e};
  \node[ex] (e7) at (6.6,-0.6) {007e};
  \node[client] (c) at (-2.4,1.2) {008c 客户端};
  \node[font=\ttfamily\scriptsize,below=1pt of c] {SOCKS 9008};
  \node[client,fill=torverde!10,draw=torverde!60!black] (h) at (2,-1.2) {009h 洋葱服务};
  \node[font=\ttfamily\scriptsize,below=1pt of h] {5858 $\to$ 127.0.0.1:4747};
  \node[server] (w) at (9.4,1.2) {\texttt{127.0.0.1:8080}};
  \begin{scope}[on background layer]
    \foreach \n in {r3,r4,r5,e6,e7} \draw[dotted,torgrigio] (\n) -- (a1);
  \end{scope}
  \draw[cellaf] (c) -- (r3);
  \draw[cellaf] (r3) to[bend right=40] (r5);
  \draw[cellaf] (r5) -- (e6);
  \draw[cellaf] (e6) -- (w);
  \node[font=\sffamily\scriptsize,text=torviola] at (2,-0.05) {HTTP 示例电路};
  \node[font=\sffamily\tiny,text=torgrigio] at (4,-2.4) {洋葱服务：无出口};
\end{tikzpicture}
